% Zone „Kennst du schon" – Nr. 1–11
\setcounter{aufgabe}{0}
\einheitenkopf[zone][Kennst du schon]{Das kennst du schon}

\begin{aufgabe}{Ich kann Zahlen quadrieren, auch negative Zahlen und Dezimalzahlen.}
Berechne.
\begin{teilezwei}
\tz $7^2 = $ \leerfeld & \tz $12^2 = $ \leerfeld \\
\tz $(-9)^2 = $ \leerfeld & \tz $0{,}6^2 = $ \leerfeld \\
\tz $3\cdot(-2)^2 = $ \leerfeld & \\
\end{teilezwei}
\end{aufgabe}

\begin{aufgabe}{Ich kann mit Klammern und negativen Zahlen rechnen.}
Berechne.
\begin{teilezwei}
\tz $4\cdot(2 + 3) = $ \leerfeld & \tz $18 - 3\cdot 4 = $ \leerfeld \\
\tz $(-2)\cdot(-2 + 6) = $ \leerfeld & \tz $(-4)\cdot(-4 - 3) = $ \leerfeld \\
\tz $2\cdot(-3)^2 - 7 = $ \leerfeld & \\
\end{teilezwei}
\end{aufgabe}

\begin{aufgabe}{Ich kann durch Einsetzen prüfen, ob eine Zahl eine Gleichung erfüllt.}
Setze die Zahl ein und kreuze an.
\begin{teile}
\teil $x = 3$ in $2x + 1 = 7$ \hfill \kreuz{(wA)} \quad \kreuz{(fA)}
\teil $x = 5$ in $x - 2 = 4$ \hfill \kreuz{(wA)} \quad \kreuz{(fA)}
\teil $x = -2$ in $3x + 10 = 4$ \hfill \kreuz{(wA)} \quad \kreuz{(fA)}
\teil $x = -2$ in $x\cdot(x + 6) = -8$ \hfill \kreuz{(wA)} \quad \kreuz{(fA)}
\teil $x = -1$ in $4 - 2x = 2$ \hfill \kreuz{(wA)} \quad \kreuz{(fA)}
\end{teile}
\end{aufgabe}

\begin{aufgabe}{Ich finde den Fehler beim Einsetzen einer negativen Zahl.}
Mia prüft, ob $x = -5$ die Gleichung $x^2 + 3x = 10$ erfüllt:
\rechnung{-5^2 + 3\cdot(-5) &= -25 - 15 \\ &= -40 \\ -40 &\neq 10 &&\text{(fA)}}
Finde den Fehler, benenne ihn und korrigiere die Rechnung.
\end{aufgabe}

\begin{aufgabe}{Ich kann eine negative Zahl in einen Term mit Quadrat einsetzen.}
Prüfe durch Einsetzen, ob die Zahl die Gleichung erfüllt.
\begin{teile}
\teil $x = -4$ in $x^2 + 2x = 8$ \hfill \kreuz{(wA)} \quad \kreuz{(fA)}
\teil $x = -2$ in $x^2 - 3x = 4$ \hfill \kreuz{(wA)} \quad \kreuz{(fA)}
\end{teile}
\end{aufgabe}

\begin{aufgabe}{Ich kann eine lineare Gleichung lösen.}
Löse die Gleichung.
\begin{gleichungsraster}[2]
\gl{x + 6 = 11} & \gl{4x = 28} \\
\gl{x + 8 = 0} & \gl{-x = 7} \\
\gl{2x - 7 = 0} & \gl[3]{3x + 4 = x - 6} \\
\end{gleichungsraster}
\end{aufgabe}

\begin{aufgabe}{Ich kann eine Quadratwurzel ziehen.}
Berechne, wenn es geht.
\begin{teilezwei}
\tz $\sqrt{49} = $ \leerfeld & \tz $\sqrt{100} = $ \leerfeld \\
\tz $\sqrt{169} = $ \leerfeld & \tz $\sqrt{0{,}25} = $ \leerfeld \\
\tz $\sqrt{-16} = $ \leerfeld & \tz $\sqrt{36 - 11} = $ \leerfeld \\
\anweisung{Rechne mit dem Taschenrechner. Runde auf zwei Stellen nach dem Komma.}
\tz $\sqrt{7} \approx $ \leerfeld & \tz $\sqrt{13{,}5} \approx $ \leerfeld \\
\end{teilezwei}
\end{aufgabe}

\begin{aufgabe}{Ich kann den Scheitelpunkt aus der Scheitelpunktform ablesen.}
Gib den Scheitelpunkt an und kreuze an, wie die Parabel geöffnet ist.
\begin{teile}
\teil $y = (x - 2)^2 + 1$ \hfill $S($\leerfeld$)$ \quad \kreuz{nach oben} \quad \kreuz{nach unten}
\teil $y = (x - 3)^2$ \hfill $S($\leerfeld$)$ \quad \kreuz{nach oben} \quad \kreuz{nach unten}
\teil $y = x^2 - 4$ \hfill $S($\leerfeld$)$ \quad \kreuz{nach oben} \quad \kreuz{nach unten}
\teil $y = (x + 1)^2 - 5$ \hfill $S($\leerfeld$)$ \quad \kreuz{nach oben} \quad \kreuz{nach unten}
\teil $y = -(x - 2)^2 + 3$ \hfill $S($\leerfeld$)$ \quad \kreuz{nach oben} \quad \kreuz{nach unten}
\end{teile}
\end{aufgabe}

\begin{aufgabe}{Ich kann Klammern ausmultiplizieren, auch mit binomischen Formeln.}
Multipliziere aus und fasse zusammen.
\begin{gleichungsraster}[2]
\gl{x\cdot(x + 4)} & \gl{3x\cdot(x - 2)} \\
\gl{(x + 5)^2} & \gl{(x - 3)^2} \\
\gl{(x + 2)\cdot(x - 6)} & \\
\end{gleichungsraster}
\end{aufgabe}

\begin{aufgabe}{Ich kann $x$ ausklammern.}
Klammere $x$ aus.
\begin{gleichungsraster}[1]
\gl{x^2 + 9x} & \gl{x^2 - 4x} \\
\gl{3x^2 + 12x} & \gl{x^2 - x} \\
\end{gleichungsraster}
\end{aufgabe}

\begin{aufgabe}{Ich kann Terme ordnen und zusammenfassen.}
Fasse zusammen und ordne nach $x^2$, $x$ und Zahl.
\begin{gleichungsraster}[2]
\gl{3 + x^2 + 2x} & \gl{5x - 1 + x^2} \\
\gl{x^2 + 4x - 3x + 7} & \gl{6 - 2x + x^2 - 9} \\
\gl{2x^2 - 6 + x - x^2} & \gl{x^2 + 2x - (3x + 4)} \\
\end{gleichungsraster}
\end{aufgabe}
